
\documentclass[aspectratio=169,10pt]{ctexbeamer}
\usefonttheme{professionalfonts}
\usetheme{default}
\usepackage{fontspec}
\IfFontExistsTF{Noto Serif CJK SC}{\setCJKmainfont{Noto Serif CJK SC}}{\setCJKmainfont{FandolSong}}
\IfFontExistsTF{Noto Sans CJK SC}{\setCJKsansfont{Noto Sans CJK SC}}{\setCJKsansfont{FandolHei}}
\IfFontExistsTF{DejaVu Sans}{\setsansfont{DejaVu Sans}}{\setsansfont{Arial}}
\usepackage{booktabs}
\usepackage{tabularx}
\usepackage{array}
\usepackage{hyperref}
\usepackage{xcolor}
\usepackage{ragged2e}
\definecolor{Accent}{HTML}{A60000}
\definecolor{Muted}{HTML}{666666}
\definecolor{LightBG}{HTML}{F7F7F7}
\definecolor{RuleBG}{HTML}{FFF7F7}
\setbeamertemplate{navigation symbols}{}
\setbeamertemplate{footline}{\hfill\scriptsize\insertframenumber\,/\,\inserttotalframenumber\hspace{1.2em}\vspace{0.7em}}
\setbeamercolor{title}{fg=black}
\setbeamercolor{frametitle}{fg=black}
\setbeamerfont{frametitle}{series=\bfseries,size=\large}
\setbeamercolor{section in toc}{fg=black}
\setbeamertemplate{itemize item}{\textcolor{Accent}{\scriptsize$\blacktriangleright$}}
\setbeamertemplate{itemize subitem}{\textcolor{Muted}{\scriptsize$\bullet$}}
\newcolumntype{Y}{>{\RaggedRight\arraybackslash}X}
\newcommand{\Tag}[1]{\textcolor{Accent}{\scriptsize\bfseries #1}}
\newcommand{\Key}[1]{\textcolor{Muted}{\scriptsize\texttt{\detokenize{#1}}}}
\newcommand{\Evidence}[1]{\textcolor{Accent}{\scriptsize\bfseries #1}}
\newcommand{\SourceState}[1]{\textcolor{Muted}{\scriptsize #1}}
\newcommand{\Meta}[4]{%
  {\Key{#1}\quad\Evidence{#2}\par
  \SourceState{#3}\par
  \SourceState{#4}\par}}
\newcommand{\Divider}{\vspace{0.6mm}\hrule\vspace{1.5mm}}
\newcommand{\RuleBox}[1]{%
  \vspace{1mm}
  \fcolorbox{Accent}{RuleBG}{\begin{minipage}{0.96\textwidth}\scriptsize #1\end{minipage}}
  \vspace{1mm}}
\newcommand{\DirectionIntro}[4]{%
\begin{frame}[t]{#1}
\small
\textbf{方向定位}\par #2
\medskip
\textbf{三篇主讲选择规则}\par #3
\medskip
\textbf{证据边界}\par #4
\end{frame}}
\newcommand{\DirectionSummary}[5]{%
\begin{frame}[t]{#1：技术演进总结}
\small
\begin{itemize}
  \item \textbf{基础问题：}#2
  \item \textbf{第一条推进线：}#3
  \item \textbf{第二条推进线：}#4
  \item \textbf{仍需保留的缺口：}#5
\end{itemize}
\end{frame}}
\newcommand{\PaperStory}[8]{%
\begin{frame}[t]{#1}
\small
#2
\Divider
\begin{itemize}
  \item 本组 slide 固定按 7 页展开：标题页、内容摘要、研究背景、核心洞察/解决方案、关键技术、实验结果/证据状态、文章评价。
  \item 证据等级在每页保留；draft、metadata-only、gated、source-limited 和 industry evidence 不写成已完全落地的学术系统。
\end{itemize}
\end{frame}
\begin{frame}[t]{#1：内容摘要}
\small
#2
\Divider
#3
\end{frame}
\begin{frame}[t]{#1：研究背景}
\small
#2
\Divider
#4
\end{frame}
\begin{frame}[t]{#1：核心洞察 / 解决方案}
\small
#2
\Divider
#5
\end{frame}
\begin{frame}[t]{#1：关键技术}
\small
#2
\Divider
#6
\end{frame}
\begin{frame}[t]{#1：实验结果 / 证据状态}
\small
#2
\Divider
#7
\end{frame}
\begin{frame}[t]{#1：文章评价}
\small
#2
\Divider
#8
\end{frame}}
\newcommand{\PaperFrame}[7]{%
\begin{frame}[t]
  \frametitle{#1}
  \vspace{-2mm}
  {\Tag{#2}}\Divider
  \begin{columns}[t,totalwidth=\textwidth]
    \begin{column}{0.49\textwidth}
      {\scriptsize\textbf{内容摘要}\par #3\par\medskip
      \textbf{研究背景}\par #4\par\medskip
      \textbf{解决方案}\par #5\par}
    \end{column}
    \begin{column}{0.49\textwidth}
      {\scriptsize\textbf{实验结果}\par #6\par\medskip
      \textbf{文章评价}\par #7\par}
    \end{column}
  \end{columns}
\end{frame}}
